\chapter{วิศวกรรมฟิสิกส์เกษตรดิจิทัล\index{ฟิสิกส์เกษตรดิจิทัล (Digital Agriphysics)}\index{Digital Agriphysics}และสรีรวิทยาพืชและสัตว์น้ำชายฝั่ง}
\label{chap:ch01}

\section*{แผนบริหารการสอนประจำบทที่ 1}
\addcontentsline{toc}{section}{แผนบริหารการสอนประจำบทที่ 1}

\begin{planbox}{หัวข้อเนื้อหาประจำบท}
\begin{enumerate}[topsep=2pt, itemsep=1pt, partopsep=0pt, parsep=0pt, leftmargin=1.5em]
    \item ความสำคัญของวิศวกรรมฟิสิกส์เกษตรดิจิทัลต่อการผลิตผลไม้มูลค่าสูงและการเพาะเลี้ยงสัตว์น้ำ
    \item รากฐานฟิสิกส์และเคมีไฟฟ้าของหัววัดสารอาหารและสภาพดินแบบบูรณาการ 7-in-1
    \item การแพร่ของก๊าซและอุณหพลศาสตร์แรงดึงระเหยน้ำของบรรยากาศสำหรับพืชสวน
    \item ฟิสิกส์การวัดออกซิเจนละลายน้ำเชิงแสงและสมดุลเคมีของน้ำในบ่อเพาะเลี้ยงสัตว์น้ำ
    \item โปรโตคอลการสื่อสารข้อมูลเซนเซอร์อุตสาหกรรมด้วย Modbus RTU\index{โปรโตคอลมอดบัส (Modbus RTU)}\index{Modbus RTU} ผ่านสายสัญญาณ RS485
\end{enumerate}
\end{planbox}

\begin{planbox}{วัตถุประสงค์เชิงพฤติกรรม}
\begin{enumerate}[topsep=2pt, itemsep=1pt, partopsep=0pt, parsep=0pt, leftmargin=1.5em]
    \item อธิบายสมการเนิร์นสต์และความสัมพันธ์เชิงฟิสิกส์ของหัววัดค่ากรด-ด่างและความเค็มได้อย่างถูกต้อง
    \item คำนวณค่าแรงดึงระเหยน้ำของอากาศจากอุณหภูมิและความชื้นสัมพัทธ์ได้อย่างแม่นยำ
    \item อธิบายหลักการดับการเรืองแสงเชิงแสงในการวัดออกซิเจนละลายน้ำได้อย่างถูกต้อง
    \item ออกแบบการเชื่อมต่อโครงข่ายเซนเซอร์ RS485 ร่วมกับการคำนวณรหัสตรวจสอบความถูกต้อง CRC-16 ได้อย่างสมบูรณ์
\end{enumerate}
\end{planbox}

\newpage

\begin{planbox}{วิธีสอนและกิจกรรมการเรียนการสอน}
\begin{enumerate}[topsep=2pt, itemsep=1pt, partopsep=0pt, parsep=0pt, leftmargin=1.5em]
    \item บรรยายหลักการทางฟิสิกส์ สรีรวิทยาพืช และสมดุลเคมีสิ่งแวดล้อมประกอบสื่อทัศนูปกรณ์
    \item สาธิตการสอบเทียบหัววัดดิน 7-in-1 และหัววัด Optical DO ด้วยสารละลายบัฟเฟอร์มาตรฐาน
    \item ปฏิบัติการทดลองวัดค่าจริงในตัวอย่างดินสวนทุเรียนและตัวอย่างน้ำจากบ่อเลี้ยงกุ้งขาว
    \item กิจกรรมกลุ่มวิเคราะห์ผลกระทบของอุณหภูมิต่อค่าคลาดเคลื่อนของเซนเซอร์และการชดเชยค่า
\end{enumerate}
\end{planbox}

\begin{planbox}{สื่อการเรียนการสอน}
\begin{enumerate}[topsep=2pt, itemsep=1pt, partopsep=0pt, parsep=0pt, leftmargin=1.5em]
    \item เอกสารคำสอนบทที่ 1 สไลด์นำเสนอ และชุดจำลองสมการฟิสิกส์
    \item ชุดหัววัดดิน 7-in-1 อุตสาหกรรม (RS485) และหัววัด Optical DO ความแม่นยำสูง
    \item บอร์ดไมโครคอนโทรลเลอร์ ESP32-S3 พร้อมโมดูลแปลงสัญญาณ MAX485
    \item สารละลายมาตรฐาน pH 4.01, 7.00, 10.01 และสารละลายมาตรฐานความเค็ม
\end{enumerate}
\end{planbox}

\begin{planbox}{การวัดและประเมินผล}
\begin{enumerate}[topsep=2pt, itemsep=1pt, partopsep=0pt, parsep=0pt, leftmargin=1.5em]
    \item การประเมินความรู้ผ่านแบบทดสอบย่อยเชิงทฤษฎีและสมการฟิสิกส์
    \item การประเมินทักษะการต่อวงจรและการสอบเทียบเซนเซอร์ในห้องปฏิบัติการ
    \item การประเมินความถูกต้องของรายงานการวิเคราะห์ชุดข้อมูลและการเขียนโปรแกรมรับส่งข้อมูล
\end{enumerate}
\end{planbox}

\clearpage

\section{ความสำคัญของฟิสิกส์เกษตรดิจิทัลในบริบทภาคตะวันออก}

การทำเกษตรกรรมแม่นยำและการเพาะเลี้ยงสัตว์น้ำในยุคดิจิทัล \textbf{ฟิสิกส์เกษตรดิจิทัล (Digital Agriphysics)} มีบทบาทเป็นสะพานเชื่อมสำคัญระหว่างปรากฏการณ์ทางกายภาพของสิ่งแวดล้อมกับระบบตัดสินใจอัตโนมัติ โดยเฉพาะในพื้นที่ภาคตะวันออกของประเทศไทย ซึ่งมีพืชเศรษฐกิจมูลค่าสูง เช่น ทุเรียนหมอนทอง (*Durio zibethinus* Murr.) และสัตว์น้ำส่งออกที่สำคัญ เช่น กุ้งขาวแวนนาไม (*Litopenaeus vannamei*) 

สภาวะทางกายภาพและเคมีของดินและน้ำส่งผลโดยตรงต่อสรีรวิทยาและการรอดชีวิตของสิ่งมีชีวิต การขาดความเข้าใจในหลักการทางวิทยาศาสตร์พื้นฐานมักนำไปสู่ความสูญเสียทางเศรษฐกิจ เช่น การตัดทุเรียนในระยะที่ไม่ได้คุณภาพเนื้อแห้ง การสูญเสียพลังงานไฟฟ้าในเครื่องตีน้ำบ่อกุ้งเกินความจำเป็น หรือการระบาดของโรครากเน่าโคนเน่าอันเนื่องมาจากระดับความชื้นในดินสะสมเกินพิกัดความจุสนาม (Field Capacity)

\section{รากฐานฟิสิกส์และเคมีไฟฟ้าของหัววัดดินบูรณาการ}

การตรวจวัดคุณสมบัติของดินตามหลักการแรกประกอบด้วยตัวแปรสำคัญทางฟิสิกส์และเคมี ได้แก่ ปริมาณธาตุอาหารหลัก ไนโตรเจน (N) ฟอสฟอรัส (P) โพแทสเซียม (K) ค่าความเป็นกรด-ด่าง (Soil pH) ค่าสภาพนำไฟฟ้า (Electrical Conductivity หรือ EC) และปริมาณความชื้นในดินเชิงปริมาตร (Volumetric Water Content หรือ VWC)

\subsection{สมการเนิร์นสต์และการชดเชยอุณหภูมิ}

การวัดค่าความเป็นกรด-ด่างในสารละลายดินอาศัยหลักการความต่างศักย์ไฟฟ้าเคมีข้ามเยื่อแก้วไวไอออนไฮโดรเจนตาม \textbf{สมการเนิร์นสต์ (Nernst Equation)\index{สมการเนิร์นสต์ (Nernst Equation)}\index{Nernst Equation}} ซึ่งอธิบายความสัมพันธ์ระหว่างแรงเคลื่อนไฟฟ้า $E$ กับกิจกรรมของไฮโดรเจนไอออนได้ดังนี้

\begin{equation}
E = E^0 - \frac{2.303 R T}{n F} \text{pH}
\label{eq:nernst}
\end{equation}

โดยที่ $E^0$ แทนศักย์ไฟฟ้ามาตรฐานของขั้วไฟฟ้า (Standard Electrode Potential), $R$ แทนค่าคงที่ก๊าซสากล ($8.314\text{ J}\cdot\text{mol}^{-1}\cdot\text{K}^{-1}$), $T$ แทนอุณหภูมิสัมบูรณ์ในหน่วยเคลวิน (K), $F$ แทนค่าคงที่ฟาราเดย์ ($96,485\text{ C}\cdot\text{mol}^{-1}$), และ $n$ แทนจำนวนประจุของไอออน ($n=1$ สำหรับ $\text{H}^+$)

พิจารณาค่าความชันเนิร์นสเตียน (Nernstian Slope) ที่อุณหภูมิต่างๆ จะพบว่ามีการเปลี่ยนแปลงตามอุณหภูมิดังตารางต่อไปนี้

\begin{table}[H]
\centering
\caption{ค่าความชันเนิร์นสเตียนตามระดับอุณหภูมิของดิน}
\label{tab:nernst_slope}
\begin{tabularx}{\textwidth}{cccY}
    \toprule
    \textbf{อุณหภูมิ ($^\circ\text{C}$)} & \textbf{อุณหภูมิสัมบูรณ์ (K)} & \textbf{ความชันตามทฤษฎี (mV/pH)} & \textbf{การชดเชยที่ pH 4.01 เทียบกับ 25 $^\circ\text{C}$ (mV)} \\
    \midrule
    20 & 293.15 & 58.16 & +2.99 \\
    25 & 298.15 & 59.16 & 0.00 (จุดอ้างอิง) \\
    30 & 303.15 & 60.15 & -2.96 \\
    35 & 308.15 & 61.14 & -5.92 \\
    40 & 313.15 & 62.13 & -8.91 \\
    \bottomrule
\end{tabularx}
\end{table}

\subsection{การคำนวณแรงดึงระเหยน้ำของอากาศสำหรับพืชสวน}

การคายน้ำของพืชถูกควบคุมโดยสภาวะทางอุณหพลศาสตร์ของอากาศรอบทรงพุ่ม \textbf{ค่าแรงดึงระเหยน้ำของอากาศ (Vapor Pressure Deficit หรือ VPD)\index{แรงดึงระเหยน้ำของอากาศ (VPD)}\index{Vapor Pressure Deficit (VPD)}} เป็นตัวชี้วัดที่แม่นยำกว่าความชื้นสัมพัทธ์ โดยคำนวณจากผลต่างระหว่างความดันไอน้ำอิ่มตัว $e_s(T)$ กับความดันไอน้ำจริง $e_a$ ดังสมการ

\begin{equation}
e_s(T) = 0.61078 \exp\left( \frac{17.27 T}{T + 237.3} \right) \quad (\text{kPa})
\label{eq:tetens}
\end{equation}

\begin{equation}
\text{VPD} = e_s(T) \times \left( 1 - \frac{\text{RH}}{100} \right) \quad (\text{kPa})
\label{eq:vpd}
\end{equation}

สำหรับต้นทุเรียนในระยะออกดอกและติดผล ช่วงค่า VPD ที่เหมาะสมสำหรับการเปิดปากใบและการสังเคราะห์ด้วยแสงจะอยู่ระหว่าง $0.8 - 1.2\text{ kPa}$ หากค่า VPD สูงเกินกว่า $1.6\text{ kPa}$ พืชจะปิดปากใบเพื่อป้องกันการสูญเสียน้ำ นำไปสู่การชะงักงันของการดูดซึมธาตุอาหารแคลเซียมและโบรอน

\section{ฟิสิกส์การวัดออกซิเจนละลายน้ำเชิงแสงในบ่อเพาะเลี้ยงสัตว์น้ำ}

ในระบบฟาร์มเพาะเลี้ยงสัตว์น้ำชายฝั่ง ออกซิเจนละลายน้ำ (Dissolved Oxygen หรือ DO) เป็นปัจจัยวิกฤติต่อการเจริญเติบโตและอัตรารอดของกุ้งขาวแวนนาไม หัววัดออกซิเจนละลายน้ำยุคใหม่ใช้หลักการ \textbf{การดับการเรืองแสงเชิงแสง (Optical Luminescence Quenching)\index{การวัดออกซิเจนละลายน้ำเชิงแสง (Optical DO)}\index{Optical Dissolved Oxygen}}

เซนเซอร์จะปล่อยแสงสีน้ำเงินความยาวคลื่น $470\text{ nm}$ ไปกระตุ้นโมเลกุลสารเรืองแสงจำพวกรูทีเนียม (Ruthenium Complex) โมเลกุลออกซิเจนอิสระที่แพร่เข้ามาจะทำปฏิกิริยาแย่งรับพลังงาน (Collisional Quenching) ส่งผลให้ความเข้มและช่วงเวลาการเรืองแสงสีแดง ($620\text{ nm}$) ลดลงตามสมการสเติร์น-โวลเมอร์ (Stern-Volmer Equation)

\begin{equation}
\frac{\tau_0}{\tau} = 1 + K_{SV} [\text{O}_2]
\label{eq:stern_volmer}
\end{equation}

โดยที่ $\tau_0$ แทนช่วงเวลาการเรืองแสงเมื่อไม่มีออกซิเจน, $\tau$ แทนช่วงเวลาการเรืองแสงที่วัดได้จริง, $K_{SV}$ แทนค่าคงที่การดับแสงของสเติร์น-โวลเมอร์, และ $[\text{O}_2]$ แทนความเข้มข้นของออกซิเจนละลายน้ำ

\begin{theorybox}[ข้อได้เปรียบเชิงฟิสิกส์ของหัววัด Optical DO เทียบกับกัลวานิก]
หัววัดเชิงแสงไม่มีการบริโภคโมเลกุลออกซิเจนในระหว่างการวัด ทำให้ไม่จำเป็นต้องพึ่งพาการไหลวนของน้ำผ่านหัววัด (Zero Flow Dependency) และไม่มีปัญหาเรื่องสารละลายอิเล็กโทรไลต์เสื่อมสภาพหรือแผ่นเมมเบรนอุดตันจากคราบเมือกแบคทีเรียในบ่อกุ้ง จึงมีเสถียรภาพการวัดระยะยาวเหนือกว่าหัววัดระบบกัลวานิกดั้งเดิมอย่างเห็นได้ชัด
\end{theorybox}

\section{โปรโตคอลการสื่อสารข้อมูลเซนเซอร์อุตสาหกรรมด้วย Modbus RTU}

การต่อเชื่อมเซนเซอร์ตรวจวัดดินและน้ำในฟาร์มจริงจำเป็นต้องส่งผ่านสายสัญญาณระยะไกลท่ามกลางสัญญาณรบกวนแม่เหล็กไฟฟ้าสูง มาตรฐานสัญญาณไฟผลต่าง \textbf{RS485 Differential Signaling} ร่วมกับโปรโตคอล \textbf{Modbus RTU} จึงเป็นมาตรฐานอุตสาหกรรมที่เหมาะสมที่สุด

\subsection{อัลกอริทึมการตรวจสอบความถูกต้องด้วยรหัสตรวจนับส่วนเกิน}

การส่งแพ็กเกตข้อมูล Modbus RTU ทุกชุดต้องลงท้ายด้วยรหัสตรวจสอบข้อผิดพลาด \textbf{CRC-16 Checksum} โดยใช้พหุนามกำเนิด $0\text{xA001}$ (Modbus Polynomial) โค้ดภาษา Python ด้านล่างแสดงการอ่านค่าเซนเซอร์ดินและตรวจสอบข้อมูล

\begin{pythonbox}[การอ่านค่าเซนเซอร์ Modbus RTU และตรวจสอบความถูกต้อง CRC-16]
def calculate_crc16(data: bytes) -> int:
    crc = 0xFFFF
    for pos in data:
        crc ^= pos
        for _ in range(8):
            if (crc & 0x0001) != 0:
                crc = (crc >> 1) ^ 0xA001
            else:
                crc >>= 1
    return crc

def parse_soil_sensor_response(response: bytes) -> dict:
    if len(response) < 19:
        raise ValueError("ขนาดแพ็กเกตข้อมูลไม่สมบูรณ์")
    payload = response[:-2]
    received_crc = response[-2] | (response[-1] << 8)
    if calculate_crc16(payload) != received_crc:
        raise ValueError("ข้อมูลผิดพลาด CRC-16 ไม่สอดคล้องกัน")
    
    # สกัดข้อมูลตามโครงสร้างรีจิสเตอร์ Modbus
    moisture = ((response[3] << 8) | response[4]) / 10.0
    temperature = ((response[5] << 8) | response[6]) / 10.0
    ec = (response[7] << 8) | response[8]
    ph = ((response[9] << 8) | response[10]) / 10.0
    nitrogen = (response[11] << 8) | response[12]
    phosphorus = (response[13] << 8) | response[14]
    potassium = (response[15] << 8) | response[16]
    
    return {
        "moist_pct": moisture,
        "temp_c": temperature,
        "ec_us_cm": ec,
        "ph": ph,
        "n_mg_kg": nitrogen,
        "p_mg_kg": phosphorus,
        "k_mg_kg": potassium
    }
\end{pythonbox}

\section*{คำถามทบทวนท้ายบทที่ 1}
\addcontentsline{toc}{section}{คำถามทบทวนท้ายบทที่ 1}

\begin{enumerate}
    \item จงอธิบายสาเหตุเชิงอุณหพลศาสตร์ที่ทำให้ความชันของหัววัดค่ากรด-ด่างตามสมการเนิร์นสต์เพิ่มขึ้นเมื่ออุณหภูมิของดินสูงขึ้น
    \item หากวัดอุณหภูมิใบพืชได้ 32 องศาเซลเซียส และความชื้นสัมพัทธ์ของอากาศเท่ากับ 55\% จงคำนวณหาค่า VPD พร้อมวิเคราะห์ว่าปากใบพืชจะอยู่ในสภาวะใด
    \item จงเปรียบเทียบข้อดีและข้อจำกัดระหว่างหัววัดออกซิเจนละลายน้ำแบบกัลวานิกกับแบบเรืองแสงเชิงแสงในบ่อเลี้ยงกุ้งที่มีความเค็ม 25 ppt
    \item เหตุใดการส่งข้อมูลผ่านสายสัญญาณ RS485 จึงทนทานต่อสัญญาณรบกวนจากมอเตอร์และปั๊มน้ำได้ดีกว่ามาตรฐาน RS232
\end{enumerate}
\clearpage
